\section{Dataset and benchmark construction}
\label{sec:dataset}

\subsection{How eval3 was built}

The evaluation set, \textbf{eval3}, was constructed so that its gold answers are fixed \emph{before any text
exists}. A skeleton generator produces groups of records about one (entity, attribute) pair --- the structure, the
record dates, the intended relation between the records and therefore the correct answer --- as structured briefs
with no prose. Only then were the record texts written, by Qwen2.5-14B-Instruct from the briefs alone. The writer
input contained no gold field, no date or source field and no label words, and the writer was required to keep each
record's own value verbatim and to write date \emph{values} in ISO form so that the mechanical validator would not
reject correct long-form dates. Questions about the current state use templates shared with every split.

A mechanical validator then rejected items whose text failed to carry its own value, whose value appeared in the
question, or that contained no entity. Of 276 skeleton items, \textbf{268 remained valid}, spanning
\textbf{135 families}. The eight exclusions are recorded by identifier.

The set was sealed on the cluster before any confirmatory run: a hash manifest over the task file, the skeleton,
the groups and the annotations, with the manifest's own hash recorded in the freeze record. The runner refuses any
task file whose hash differs from the sealed one.

\subsection{Annotation, and why it is a sensitivity analysis rather than a filter}

eval3 has \textbf{no human validation}. It was annotated blind by two model annotators against the three gold
classes. Agreement with gold was 174 of 268 ($\kappa=0.471$) for one annotator and 232 of 268 ($\kappa=0.767$) for
the other \citep{cohen1960kappa}; 161 items have double agreement and 107 are flagged by at least one annotator.

These numbers are annotation statistics, not outcomes, and they are reported here because they bound what the set
can support. We do \emph{not} filter on them: the primary analysis uses all 268 valid items, and the
double-agreement subset (161), the flagged complement (107) and the subset agreed by the stronger annotator (232)
are three separate sensitivity analyses (\Cref{tab:sensitivity}). Effects are somewhat \emph{larger} on the
flagged complement ($+21.3$\pp{} on H-primary) than on the double-agreement subset ($+14.3$\pp), which is the
opposite of what an ``annotator noise inflates the effect'' account would predict.

\subsection{What the set is, and is not}

eval3 is synthetic, templated, English-only, ISO-dated, model-written and model-annotated. It was built to be
sealed and pre-registered against, and that is the only claim we make for it. It is not a benchmark of record: the
closest external set, STALE, has 400 expert-validated scenarios, and DISCERN's evaluation set is human-vetted
\citep{chao2026stale,ahuja2026discern}. One consequence is worth stating plainly: the writer, Qwen2.5-14B, is also
one of the annotators and one of the evaluated secondary models, so that model's results carry a dependence the
others do not. We report it separately and never pool it.

A second consequence runs the other way. Because eval3 \emph{includes} explicitly signalled invalidation --- the
corrections, supersessions and negations a lexical arbiter is built to read --- it is favourable ground for Evarb
by construction. STALE deliberately excludes those cases, which makes it the natural external test of what this
mechanism cannot do, and an obvious next step rather than a result we can claim here.
